% ============================================================================
%  Ustawienia dokumentu (preambuła)
% ============================================================================

% --- Język i czcionki --------------------------------------------------------
\usepackage[T1]{fontenc}
\usepackage[utf8]{inputenc}
\usepackage[english,polish]{babel}
\usepackage{amsmath}
\usepackage{newtxtext,newtxmath}
\usepackage[scaled=0.84]{beramono}
\usepackage{microtype}

% --- Układ strony -------------------------------------------------------------
\usepackage[a4paper,top=2.5cm,bottom=2.5cm,inner=3.0cm,outer=2.0cm,footskip=1.1cm]{geometry}
\usepackage{setspace}
\setstretch{1.25}
\setlength{\parindent}{1.25cm}
\setlength{\parskip}{0pt plus 1pt}
\widowpenalty=10000
\clubpenalty=10000
\raggedbottom

% --- Pakiety pomocnicze ------------------------------------------------------
\usepackage{graphicx}
\usepackage{xcolor}
\usepackage{float}
\usepackage{array}
\usepackage{booktabs}
\usepackage{tabularx}
\usepackage{multirow}
\usepackage{longtable}
\usepackage{enumitem}
\usepackage{etoolbox}
\usepackage{dirtree}
\usepackage{multicol}
\usepackage{url}
\usepackage{tikz}
\usetikzlibrary{positioning,arrows.meta,shapes.geometric,shapes.multipart,fit,calc,backgrounds}

% --- Nagłówki rozdziałów -----------------------------------------------------
\usepackage[explicit]{titlesec}
\titleformat{\chapter}[block]
  {\normalfont\fontsize{16}{20}\selectfont\bfseries}{}{0pt}
  {\ifnum\value{chapter}>0 \thechapter.\hspace{0.5em}\fi#1}
\titleformat{name=\chapter,numberless}[block]
  {\normalfont\fontsize{16}{20}\selectfont\bfseries}{}{0pt}{#1}
\titlespacing*{\chapter}{0pt}{0pt}{16pt}
\titleformat{\section}
  {\normalfont\fontsize{13}{16}\selectfont\bfseries}{\thesection.}{0.5em}{#1}
\titlespacing*{\section}{0pt}{11pt}{4pt}
\titleformat{\subsection}
  {\normalfont\fontsize{12}{15}\selectfont\bfseries}{\thesubsection.}{0.5em}{#1}
\titlespacing*{\subsection}{0.6cm}{8pt}{3pt}

% --- Spis treści -------------------------------------------------------------
\usepackage[titles]{tocloft}
\renewcommand{\cfttoctitlefont}{\normalfont\fontsize{16}{20}\selectfont\bfseries}
\renewcommand{\cftaftertoctitle}{\par\vspace{6pt}}
\setlength{\cftbeforetoctitleskip}{0pt}
\setlength{\cftaftertoctitleskip}{8pt}
\renewcommand{\cftchapfont}{\normalfont}
\renewcommand{\cftchappagefont}{\normalfont}
\renewcommand{\cftchapleader}{\cftdotfill{\cftdotsep}}
\renewcommand{\cftchapaftersnum}{.}
\renewcommand{\cftsecaftersnum}{.}
\renewcommand{\cftsubsecaftersnum}{.}
\setlength{\cftbeforechapskip}{5pt}
\setlength{\cftchapindent}{0.4cm}
\setlength{\cftsecindent}{1.2cm}
\setlength{\cftsubsecindent}{2.2cm}
\setlength{\cftchapnumwidth}{0.8cm}
\setlength{\cftsecnumwidth}{1.0cm}
\setlength{\cftsubsecnumwidth}{1.3cm}
\setcounter{tocdepth}{1}

% --- Stopka z numerem strony (strony nieparzyste po prawej, parzyste po lewej)
\usepackage{fancyhdr}
\fancyhf{}
\fancyfoot[LE,RO]{\thepage}
\renewcommand{\headrulewidth}{0pt}
\fancypagestyle{plain}{\fancyhf{}\fancyfoot[LE,RO]{\thepage}\renewcommand{\headrulewidth}{0pt}}

% --- Podpisy rysunków, tabel i listingów -------------------------------------
\usepackage{caption}
\captionsetup{labelsep=period,font={small,stretch=1.05},justification=centering,
  singlelinecheck=false,skip=6pt}
\captionsetup[table]{position=above}
\captionsetup[figure]{position=below}

% --- Listingi kodu C# ---------------------------------------------------------
\usepackage{listings}
\definecolor{kodSlowo}{RGB}{20,40,130}
\definecolor{kodKomentarz}{RGB}{105,105,105}
\definecolor{kodTekst}{RGB}{140,30,30}
\definecolor{kodRamka}{RGB}{170,170,170}
\renewcommand{\lstlistingname}{Listing}
\renewcommand{\lstlistlistingname}{Spis listingów}
\lstdefinestyle{csharp}{
  language=[Sharp]C,
  basicstyle=\linespread{1.0}\ttfamily\footnotesize,
  keywordstyle=\bfseries\color{kodSlowo},
  commentstyle=\itshape\color{kodKomentarz},
  stringstyle=\color{kodTekst},
  morekeywords={var,async,await,readonly,get,set,nameof,where},
  showstringspaces=false,
  columns=fullflexible,
  keepspaces=true,
  breaklines=true,
  breakatwhitespace=true,
  tabsize=4,
  xleftmargin=0.6cm,
  xrightmargin=0.2cm,
  frame=tb,
  framerule=0.4pt,
  rulecolor=\color{kodRamka},
  framesep=4pt,
  aboveskip=10pt,
  belowskip=2pt,
  captionpos=b,
  literate={ą}{{\k{a}}}1 {ę}{{\k{e}}}1 {ó}{{\'o}}1 {ś}{{\'s}}1 {ł}{{\l{}}}1
           {ż}{{\.z}}1 {ź}{{\'z}}1 {ć}{{\'c}}1 {ń}{{\'n}}1
           {Ą}{{\k{A}}}1 {Ę}{{\k{E}}}1 {Ó}{{\'O}}1 {Ś}{{\'S}}1 {Ł}{{\L{}}}1
           {Ż}{{\.Z}}1 {Ź}{{\'Z}}1 {Ć}{{\'C}}1 {Ń}{{\'N}}1
}
\lstdefinestyle{json}{
  basicstyle=\linespread{1.0}\ttfamily\scriptsize,
  showstringspaces=false,
  columns=fullflexible,
  keepspaces=true,
  frame=none,
  xleftmargin=0.2cm,
  aboveskip=2pt,
  belowskip=2pt,
  stringstyle=\color{kodTekst},
  morestring=[b]",
}
\lstset{style=csharp}

% Tabele i drzewa katalogów składane z mniejszą interlinią.
\AtBeginEnvironment{table}{\small\setstretch{1.08}}
\AtBeginEnvironment{tabularx}{\setstretch{1.08}}
\setlength{\DTbaselineskip}{13pt}
\renewcommand*\DTstylecomment{\rmfamily\itshape\small}
\renewcommand*\DTstyle{\ttfamily\small}

% --- Odnośniki ---------------------------------------------------------------
\usepackage[hidelinks,unicode,bookmarksnumbered]{hyperref}
\urlstyle{same}
\def\UrlBreaks{\do\/\do\-\do\_\do\.\do\?\do\=\do\&}
% Ścieżki i identyfikatory z kropkami, łamane w miejscu "/" lub ".".
\DeclareUrlCommand\sciezka{\def\UrlFont{\ttfamily\small}}
\setlength{\emergencystretch}{3em}

% --- Wpisy w spisach rysunków, tabel i listingów -----------------------------
% Format zgodny z wcześniejszymi pracami: "Rysunek 1.1. Podpis; Źródło: ... ..... 4".
\makeatletter
\newcommand{\spisWpis}[3]{%
  \begingroup
  \def\numberline##1{#1~##1.\ }%
  \@dottedtocline{1}{0em}{0em}{#2}{#3}%
  \endgroup}
\AtBeginDocument{%
  \renewcommand*{\l@figure}[2]{\spisWpis{Rysunek}{#1}{#2}}%
  \renewcommand*{\l@table}[2]{\spisWpis{Tabela}{#1}{#2}}%
  \renewcommand*{\l@lstlisting}[2]{\spisWpis{Listing}{#1}{#2}}%
}
\makeatother

% --- Polecenia pomocnicze ----------------------------------------------------
\newcommand{\zrodloWlasne}{Źródło: Opracowanie własne}
\newcommand{\kod}[1]{\texttt{\small #1}}
\newcommand{\klasa}[1]{\texttt{\small #1}}

% Zrzut ekranu z gry; jeśli pliku brak, w jego miejscu pojawia się ramka.
\newcommand{\zrzut}[2][\textwidth]{%
  \IfFileExists{rysunki/gra/#2.png}%
    {\includegraphics[width=#1]{rysunki/gra/#2.png}}%
    {\fbox{\parbox[c][0.3\textheight][c]{#1}{\centering Brak pliku rysunki/gra/#2.png}}}}

% Style diagramów TikZ.
\tikzset{
  diag/.style={font=\linespread{1.0}\selectfont\footnotesize, >=Stealth},
  blok/.style={draw, rounded corners=2pt, align=center, inner sep=4pt, minimum height=8mm, fill=white},
  modul/.style={blok, fill=black!4},
  akcent/.style={blok, fill=blue!6},
  strz/.style={-{Stealth[length=2.2mm]}, thick},
  strzp/.style={-{Stealth[length=2.2mm]}, thick, dashed},
  przypadek/.style={draw, ellipse, align=center, inner sep=2pt, minimum height=9mm, fill=white},
}
\newcommand{\aktor}[2]{% nazwa węzła, podpis
  \begin{scope}[shift={(#1)}]
    \draw[thick] (0,0.55) circle (0.16);
    \draw[thick] (0,0.39) -- (0,-0.05);
    \draw[thick] (-0.28,0.25) -- (0.28,0.25);
    \draw[thick] (0,-0.05) -- (-0.22,-0.42);
    \draw[thick] (0,-0.05) -- (0.22,-0.42);
    \node[font=\footnotesize] at (0,-0.68) {#2};
  \end{scope}}
